\subsection{Groups}
\label{subsec:groups}

\begin{tcolorbox}[title={Definition --- Group}]
\begin{definition}[Group]
\label{def:group}
A \textbf{group} is a pair $(G,\star)$ where $G$ is a set and $\star$ is a
binary operation on $G$ satisfying:
\begin{enumerate}[label=\textbf{G\arabic*.}]
  \item associativity;
  \item existence of an identity element $e\in G$;
  \item existence, for each $a\in G$, of an inverse element $a^{-1}\in G$.
\end{enumerate}
\end{definition}
\end{tcolorbox}

\begin{remark*}[Standard quantified statement]
\[
\operatorname{Group}(G,\star)
\Longleftrightarrow
\exists e\in G,\;
\operatorname{Monoid}(G,\star,e)
\land
\forall a\in G,\;\exists b\in G\;
\operatorname{InverseElement}(b,a,\star,e,G).
\]
\end{remark*}

\begin{remark*}[Definition predicate reading]
$(G,\star)$ is a group exactly when it is an associative closed operation with
an identity and an inverse for every element.
\end{remark*}

\begin{remark*}[Negated quantified statement]
\[
\neg\operatorname{Group}(G,\star)
\Longleftrightarrow
\neg\operatorname{Associative}(\star,G)
\lor
\text{no two-sided identity exists}
\lor
\text{some element has no two-sided inverse}.
\]
\end{remark*}

\begin{remark*}[Interpretation]
Groups are algebraic structures of reversible composition.
\end{remark*}

\begin{remark*}[Dependencies]
\begin{itemize}
  \item Binary operation.
  \item Associativity.
  \item Identity element.
  \item Inverse element.
\end{itemize}
\end{remark*}

\begin{definition}[Order of a Group]
\label{def:order-of-a-group}
The \textbf{order} of a group $G$, denoted $|G|$, is the cardinality of its
underlying set.  If $|G|$ is finite, $G$ is a finite group; otherwise it is an
infinite group.
\end{definition}

\begin{remark*}[Standard quantified statement]
\[
\operatorname{OrderOfGroup}(G)=|G|.
\]
\end{remark*}

\begin{remark*}[Definition predicate reading]
The order of $G$ is the number of elements in the carrier set of $G$.
\end{remark*}

% Omitted Negated quantified statement: this definition introduces notation for
% a cardinal value rather than a predicate with a useful failure condition.

\begin{remark*}[Interpretation]
Finite groups are classified by a natural-number order; infinite groups have
infinite cardinal order.
\end{remark*}

\begin{remark*}[Dependencies]
\begin{itemize}
  \item Group.
  \item Cardinality.
\end{itemize}
\end{remark*}

\begin{remark*}[Examples]
Examples include $(\mathbb{Z},+)$, $(\mathbb{Q}\setminus\{0\},\cdot)$,
$(\mathbb{Z}/n\mathbb{Z},+)$, and $GL_n(\mathbb{R})$ under matrix
multiplication.
\end{remark*}
